% ===========================================================================
% 4 Experimental Setup
% Drafted with AI assistance (Claude). Rewrite in your own words before submitting.
% Revised on 4 Nov 2026 after the supervisor's review: third person, shorter.
% Settings: configs/scifact.yaml; dataset statistics computed from the BEIR copies.
% ===========================================================================
\section{Experimental Setup}
\label{sec:setup}

\subsection{Datasets}
\label{sec:datasets}

Both datasets come from the BEIR benchmark \parencite{thakur2021beir}, which provides a corpus,
test queries and relevance judgements in one format (Table~\ref{tab:datasets}).

\paragraph{SciFact} \parencite{wadden2020scifact} is the main dataset. Its corpus holds 5,183
abstracts of biomedical papers; each of its 300 test claims was written by an expert, and the
abstracts that support or refute it are judged relevant (339 in total, 1.13 per claim, all with
grade 1). Its 809 training claims are disjoint from the test claims and are used only to
fine-tune the encoder and to select the fusion weight, with 120 held out as development claims
(Section~\ref{sec:finetune}).

\paragraph{NFCorpus} \parencite{boteva2016nfcorpus} tests whether the improvement transfers
(RQ3). Its corpus holds 3,633 medical abstracts, and its queries come from NutritionFacts.org, a
website that explains nutrition research in plain language. It is a demanding test for two
reasons: each of its 323 test queries has many judged abstracts (38.2 on average, graded 1 for
indirectly and 2 for directly relevant), and its queries are short and phrased very differently
from the abstracts. Its 2,590 training queries are used in the same way as SciFact's training
claims, with 120 held out for development. Because the training queries have 42.7 judged
abstracts each, all graded 1, training uses the first three judged abstracts of each query
(6,917 pairs).

\begin{table}[t]
\centering\small
\caption{The two datasets. Passages are built as in Section~\ref{sec:chunking}.}
\label{tab:datasets}
\begin{tabular}{@{}lrr@{}}
\toprule
& SciFact & NFCorpus \\
\midrule
Abstracts in the corpus & 5,183 & 3,633 \\
Passages (3-sentence windows) & 21,015 & 16,574 \\
Test queries & 300 & 323 \\
Relevant abstracts per test query & 1.13 & 38.2 \\
Relevance grades & 1 & 1, 2 \\
Training queries (train / development) & 809 (689 / 120) & 2,590 (2,470 / 120) \\
\bottomrule
\end{tabular}
\end{table}

\subsection{Models and Settings}
\label{sec:models}

All models are public and are obtained from the Hugging Face Hub; Table~\ref{tab:settings} lists
every setting. Two configurations are compared throughout:

\begin{itemize}
\item the \textbf{standard hybrid pipeline}: the pretrained MiniLM encoder, unweighted RRF
($w = 1$) and the MS~MARCO cross-encoder reranking the top 50, evaluated in all four variants
(BM25, dense, fusion, fusion with reranking);
\item the \textbf{proposed retriever}: the MiniLM encoder fine-tuned on the dataset's training
claims, fused with BM25 at the weight selected on the development claims, without a reranker.
\end{itemize}

\noindent Further configurations change one component each: whole abstracts instead of
passages, score fusion instead of RRF, a stronger off-the-shelf encoder
(\texttt{bge-small-en-v1.5}, 33.4M parameters, \citeauthor{xiao2024cpack},
\citeyear{xiao2024cpack}) with a larger reranker (\texttt{bge-reranker-base}, 278M), the
reranking depth $M \in \{10, 25, 50, 100\}$ and the RRF constant $k \in \{10, 30, 60, 100\}$.

\begin{table}[t]
\centering\small
\caption{Settings of every stage (\texttt{configs/scifact.yaml}).}
\label{tab:settings}
\begin{tabularx}{\textwidth}{@{}lL@{}}
\toprule
Stage & Setting \\
\midrule
Passages & 3 sentences, stride 2, title prepended, windows under 120 characters end the abstract \\
BM25 & $k_1 = 0.9$, $b = 0.4$; lower-case, 52 stop words, Porter stemming; top 100 \\
Dense encoder & \texttt{all-MiniLM-L6-v2} (22.7M parameters, 384 dimensions), mean pooling, unit length, 256 tokens; exact inner-product search; top 100 \\
Fine-tuning & multiple-negatives ranking loss, scale 20, batch 32, AdamW, learning rate $2 \times 10^{-5}$, weight decay 0.01, 10\% warm-up then linear decay, 12 epochs, best development epoch kept \\
Fusion & RRF, $k = 60$; dense weight 1; BM25 weight 1 (standard) or selected on development claims from $\{0, 0.1, 0.2, 0.3, 0.5, 0.7, 1, 1.5, 2\}$; top 100 \\
Reranker & \texttt{ms-marco-MiniLM-L-6-v2} (22.7M), 512 tokens, top 50 rescored \\
Evaluation & nDCG@10 (main), Recall@$k$, MRR@10; paired bootstrap, 10,000 resamples, seed 42 \\
\bottomrule
\end{tabularx}
\end{table}

\subsection{Hardware and Reproducibility}
\label{sec:hardware}

All experiments run on one laptop with an Intel Core i5-12450H processor, 16~GB of memory and an
NVIDIA RTX~3050 Laptop GPU with 4~GB, using PyTorch 2.6, FAISS 1.15, Transformers 5.17 and
sentence-transformers 6.1. Latency is the wall-clock time per query until its final ranked list,
measured one query at a time after an untimed warm-up query and averaged over all test queries;
timings varied by up to a factor of two between identical runs, so they are compared only when
they differ by more than that. Every setting is kept in one configuration file, every random
choice uses seed 42, and every metric is computed from its definition and checked by unit tests
against values worked out by hand. A clean installation of the code reproduced all reported
retrieval metrics exactly.
